\documentclass[10pt]{article}
\usepackage[utf8]{inputenc}
\usepackage[T1]{fontenc}
\usepackage{lmodern}
\usepackage[margin=1in]{geometry}
\usepackage{graphicx}
\usepackage{booktabs}
\usepackage{caption}
\usepackage{subcaption}
\usepackage{float}
\usepackage{pdflscape}
\usepackage{array}
\usepackage{tabularx}
\usepackage{makecell}
\usepackage{microtype}
\usepackage[hidelinks]{hyperref}
\usepackage[numbers,sort&compress,round]{natbib}
\usepackage{authblk}

\title{Cross-Domain Reactivation of Collective Attention in Music and Film}

\author[1,2,3,*]{Cristian Candia}
\author[1,5]{Camila Utreras-Cifuentes}
\author[1,4]{Francisca Droguett}

\affil[1]{Computational Research in Social Science Laboratory, School of Engineering, Universidad del Desarrollo, Las Condes, Chile}
\affil[2]{Instituto de Data Science, Facultad de Ingenier\'ia, Universidad del Desarrollo, Las Condes, 7610658, Chile}
\affil[3]{Northwestern Institute on Complex Systems (NICO), Northwestern University, Evanston, IL 60208, USA}
\affil[4]{University of Chile, Santiago, Chile}
\affil[5]{Center for Research in Social Complexity (CICS), Faculty of Government, Universidad del Desarrollo, Las Condes, Chile}
\affil[*]{Corresponding Author: crcandiav@gmail.com}

\date{}

\begin{document}

\maketitle

\begin{abstract}
A song's appearance in a film can coincide with renewed attention, yet soundtracks may also select songs that already attract larger audiences. We examine these possibilities by linking 32,124 Billboard Hot 100 song-artist records, including 4,683 with observed film appearances, to four Spotify popularity snapshots, Last.fm cumulative listener reach, and film metadata. Film-linked songs have higher Spotify popularity and cumulative Last.fm reach, and broad Spotify attention differences are already present before observed placement. To assess subsequent renewal, we estimate age- and snapshot-specific expected-attention baselines from songs with no observed film appearance. Excess attention is observed Spotify popularity minus this fitted value. Prior-attention matching substantially reduces the broad excess-attention difference. Our primary temporal comparison focuses on 119 songs in the lower tercile of pre-film excess attention among 355 film-linked songs with temporal coverage. Their pre-to-post increase exceeds that of matched controls assigned comparable pseudo-event windows by 5.04 excess-attention points (integrated 95\% CI [2.97, 6.87]). Their probability of exceeding the pre-film tercile threshold is 18.0 percentage points higher (95\% CI [6.7, 27.6]). Other Billboard songs by the same artists also show pre-existing attention advantages, but matched comparisons do not establish corresponding increases in those songs. These results distinguish prior visibility from subsequent renewal and support an observational reactivation signature around film placement.
\end{abstract}

\section*{Significance Statement}

When a familiar song appears in a film, does its renewed visibility reflect attention it already commanded or a subsequent increase? Across previously charting songs, film-linked music has substantial attention advantages before placement. A temporal comparison finds greater increases among songs with relatively lower pre-film excess attention, measured against an age-specific baseline, after matching on prior popularity and chart history. Comparable increases are not established for other Billboard songs by the same artists. Distinguishing these patterns helps explain how cultural reuse relates to renewed public attention without treating existing popularity as evidence of revival or inferring how much listeners remember.

\noindent\textbf{Keywords:} collective memory, collective attention, cultural reuse, music, film

\medskip

A song can leave the charts and later enter a film's soundtrack. Its appearance in a new narrative offers an occasion for audiences to encounter it again, but a popular soundtrack song may also have remained popular all along. Distinguishing these possibilities is essential to understanding how cultural reuse relates to renewed collective attention. Music and film make the problem observable: song histories establish earlier public visibility, soundtrack records locate reuse in another medium, and listening traces allow comparisons before and after that reuse.

Collective-memory research provides a framework for understanding why cultural objects persist beyond their initial audiences \citep{Halbwachs1992,Nora1989,Assmann1995,Hirst2018}. It distinguishes the social resources and practices through which communities encounter the past from the memories individuals form through those encounters \citep{Hirst2008}. Archives and repertoires preserve material that need not remain continuously prominent in public life \citep{Assmann2008,Assmann2011}. Digital traces make some forms of revisiting this material observable \citep{yasseri2022collective}. They allow the study of attention to cultural objects, while leaving what people remember a separate empirical question.

Computational studies identify several temporal patterns relevant to this distinction \citep{Candia2025TemporalDynamicsCSS}. Research on songs and other cultural objects documents age-dependent declines in attention \citep{Candia2019}, while models of patent citations distinguish aging from cumulative advantage \citep{Higham2017}. Across several cultural domains, popular topics also show increasingly rapid turnover \citep{Lorenz-Spreen2019}. These are different processes: older objects can lose attention even without an acceleration in how quickly topics replace one another. Selective persistence as knowledge expands has also been documented in science and technology \citep{CandiaUzzi2021SelectiveForgetting}. Renewal introduces another temporal pattern. Recent aircraft accidents can draw Wikipedia attention to earlier related accidents \citep{Garcia-Gavilanes2017}, and public figures' deaths generate surges followed by declines in news and social-media attention \citep{West2021}. Together, these findings motivate evaluating a song's current attention against both its age and the attention it attracted before reuse.

Popular music connects personal recollection with shared repertoires through public practices of listening and remembering \citep{van2006record}. Research on Cultural Revolution songs in China, for example, documents generational differences in how a shared repertoire is remembered \citep{bryant2005music}. Film placement can bring a song into a new narrative, connecting musical remembrance to the circulation of cultural objects across media \citep{JenkinsDeuze2008,bolter1996remediation,erll2016travelling}. Music can evoke affect and autobiographical associations \citep{janata2007characterisation,krumhansl2013cascading}, while narrative context can organize attention and emotional engagement \citep{Green2004,BezdekGerrig2017}. These processes offer plausible reasons for a soundtrack appearance to draw listeners toward a song. They also offer reasons for filmmakers to choose songs that audiences already know. Soundtrack selection and subsequent attention renewal can therefore coexist. Our listening and linkage data can distinguish some of their observable patterns, although they do not measure recognition, nostalgia, or the reasons for a placement.

We study Billboard Hot 100 song-artist records, a population that has already achieved chart visibility. The analysis consequently concerns attention after initial success. Spotify popularity snapshots measure current attention on that platform, whereas the July 2017 Last.fm snapshot measures cumulative listener reach. Neither is a direct measure of memory. We use a collective-memory decay model to estimate an expected-attention baseline for songs of the same age within each Spotify snapshot. The resulting curves summarize cross-sectional age--attention profiles. We call the difference between observed popularity and the fitted value \emph{excess attention}. A positive excess-attention level can reflect persistent popularity or selection. We reserve \emph{reactivation signature} for a greater pre-to-post increase than that of matched comparison songs, an observational pattern consistent with renewed attention around film placement.

This distinction is particularly relevant for songs with relatively lower attention before reuse, where a later increase could also arise through regression to the mean. Our primary temporal comparison selects the lower tercile of pre-film excess attention among film-linked songs with measurements before and after their first observed film year. This is a relative classification within the eligible group. It does not imply that these songs were forgotten, had low absolute popularity, or fell below the age-specific baseline. Comparisons with songs assigned similar pseudo-event windows and matched on prior attention and chart history test whether subsequent increases exceed those seen among similarly selected controls.

Attention could also extend to other music by the same artist. Associative accounts of memory suggest one possible route \citep{collins1975spreading,Hirst2018}, and evidence on music sales shows that a new album can increase sales of an artist's earlier releases \citep{HendricksSorensen2009}. That evidence motivates an empirical question about the reach of attention following reuse in another medium. Here, \emph{artist catalog} refers specifically to other Billboard songs credited to the same normalized artist or band, rather than the artist's entire repertoire. We ask whether those songs show corresponding changes in attention.

The analyses address three questions. Do film-linked songs and their artists' other songs have attention advantages before observed placement? Do film-linked songs subsequently gain attention relative to an age-specific baseline and comparison songs with similar prior histories? Do corresponding gains appear in other Billboard songs by the same artists? Broad comparisons establish the scale of the associations, while temporal and matched comparisons address prior differences and subsequent change. Secondary analyses describe how attention varies with the number of film appearances and recorded worldwide box office. These measures characterize reuse and film-level visibility without measuring individual exposure to a soundtrack.

\section*{Results}

\subsection*{Data coverage and analytic universe}

The analytical universe contains 32,124 Billboard Hot 100 song-artist records, of which 4,683 are linked to at least one observed film appearance. The main Spotify analyses use four high-coverage current-popularity snapshots, providing 83,785 nonmissing song-snapshot observations. Expected-attention analyses use 83,780 observations with valid nonnegative song age, covering 28,310 distinct Billboard song-artist records. Last.fm provides one cumulative listener snapshot, July 2017, with 17,988 records. These broad samples are distinct from the smaller subsets with the temporal coverage and matching support required for comparisons around film placement.

For models involving film visibility, expanded movie metadata are available for 4,682 film-linked records, and worldwide box-office metadata are available for 4,272. Budget metadata are available for 2,581 film-linked records. Detailed platform coverage, linkage procedures, movie-metadata coverage, and missingness are reported in the SI Appendix, Data Construction and Platform Coverage.

\begin{table}[!htbp]
\centering
\small
\caption{Analytic data components.}
\label{tab:data-audit}
\begin{tabular}{@{}>{\raggedright\arraybackslash}p{0.27\linewidth}>{\raggedright\arraybackslash}p{0.11\linewidth}>{\raggedright\arraybackslash}p{0.56\linewidth}@{}}
\toprule
Data component & Support & Definition / description \\
\midrule
Billboard universe & 32,124 & Billboard Hot 100 song-artist records rebuilt from the weekly chart file \\
Film-linked Billboard records & 4,683 & Records linked to at least one IMDb soundtrack appearance by exact normalized song/artist keys \\
High-coverage Spotify panel & 83,785 & Song-snapshot observations across October 2016, July 2017, August 2022, and August 2025 \\
Expected-attention support & 83,780 & Song-snapshot observations with valid nonnegative song age \\
Last.fm snapshot & 17,988 & July 2017 cumulative listener reach (raw Listeners field, analyzed as log1p listeners) \\
Movie metadata & 4,682 & Film-linked records with exact internal-ID movie metadata \\
Worldwide box office & 4,272 & Film-linked records with worldwide box-office metadata \\
\bottomrule
\end{tabular}
\begin{flushleft}\footnotesize
The main Spotify analyses use only the four high-coverage snapshots listed above. Platform coverage, linkage, and movie-metadata missingness diagnostics are reported in the SI Appendix, Data Construction and Platform Coverage.
\end{flushleft}
\end{table}




\subsection*{Film-linked songs have attention advantages that precede placement}

Film-linked Billboard songs have higher Spotify popularity and cumulative Last.fm reach than songs with no observed film appearance. Adjusted Spotify attention premiums are positive in every high-coverage snapshot: 12.28 points in October 2016, 13.36 in July 2017, 14.29 in August 2022, and 15.09 in August 2025. The pooled Spotify premium across high-coverage snapshots is 13.79 points after adjusting for song age, squared age, Billboard weeks, peak position, and snapshot fixed effects. The July 2017 Last.fm snapshot shows a 1.78 log-listener premium. Figure~\ref{fig:attention-matching}A standardizes each platform/snapshot outcome for visual comparison. Raw-unit and standardized estimates are reported in the SI Appendix, Platform Coverage and Attention Differences.

These contemporary differences alone cannot distinguish persistent popularity from renewed attention. Broad Spotify event-time comparisons show that an attention advantage is already present before the first observed film appearance. Songs measured at least two years before placement are 9.50 popularity points above the omitted reference, which combines songs with no observed film appearance and observations in the immediate pre-film year. The event-year contrast is 11.23 points, the one-to-five-year post contrast is 10.16 points, and the more-than-five-year post contrast is 14.18 points. These are adjusted differences between event-time groups, not changes traced within a fixed song cohort. The positive pre-event contrast establishes prior differences consistent with selection into film use, without identifying why a song was chosen.

Matching on observed histories substantially reduces the attention advantage. Nearest-neighbor matching with replacement on chart year, song age, Billboard weeks, peak position, and artist-level Billboard success yields differences of 10.17, 11.18, 11.26, and 11.68 Spotify points across the four snapshots. After October 2016, matching can additionally balance popularity in the previous high-coverage snapshot with a five-point caliper. The corresponding differences are 1.95 points in July 2017, 2.23 in August 2022, and 1.96 in August 2025, with maximum absolute standardized mean differences below 0.05 (Figure~\ref{fig:attention-matching}B--C). Prior Spotify attention thus accounts for much of the broad association. Because the prior snapshot is not necessarily before film placement in this design, these matched level differences remain descriptive comparisons. The temporal analyses below require measurements before placement.

\begin{figure}[!htbp]
\centering
\includegraphics[width=\linewidth]{Figures/Main/figure1_attention_regimes_matching.pdf}
\caption{Film-linked songs have higher Spotify popularity and cumulative Last.fm reach. Panel A: regression-adjusted differences standardized within each platform/snapshot, with HC1 95\% CIs. Panel B: nearest-neighbor matched Spotify mean differences between film-linked songs and songs with no observed film appearance, with paired 95\% CIs. Hist.\ denotes matching on chart history and artist-level Billboard success. Prior additionally matches popularity in the previous Spotify snapshot with a five-point caliper, which need not precede film placement. Panel C: maximum absolute standardized mean differences before and after matching. The unit is the Billboard song-artist record within a platform/snapshot. Raw and standardized estimates appear in the SI Appendix section \textit{Platform Coverage and Attention Differences}. Covariate-level balance and replacement-aware inference appear in \textit{Matched Song-Level Comparisons}.}
\label{fig:attention-matching}
\end{figure}



\subsection*{Temporal comparisons show renewal beyond the expected-attention baseline}

To evaluate attention relative to song age, we fit snapshot-specific canonical collective-memory decay curves to Spotify age means for songs with no observed film appearance \citep{Candia2019,CandiaUzzi2021SelectiveForgetting}. The resulting expected-attention baseline is specific to age and snapshot. Excess attention is observed Spotify popularity minus this fitted value. The curves summarize cross-sectional age--attention profiles within snapshots, rather than observed decay paths of individual songs. Spotify popularity remains on its original 0--100 scale. Parameters, alternative models, and fit diagnostics are reported in the SI Appendix, Expected-Attention Baselines and MemoryDecay Diagnostics.

The pooled excess-attention premium for film-linked songs is 19.12 Spotify points. Restricting film exposure to appearances released by each snapshot year also yields a positive difference. These level comparisons can combine persistent popularity, selection, and subsequent increases. After matching on prior attention, the pooled difference falls to 2.52 excess-attention points over 12,095 matched pairs (replacement-aware two-way cluster-bootstrap 95\% CI [2.16, 2.87]). A positive residual difference motivates a temporal comparison, but does not itself establish renewal.

The primary temporal analysis examines songs with lower pre-film excess attention. Of 355 film-linked songs with Spotify measurements before and at or after their first observed film year, 119 fall in the bottom tercile of pre-film excess attention. This classification is relative to the eligible group and does not require low absolute popularity or attention below the fitted baseline. To address regression to the mean, we compare these songs with controls that meet the same baseline threshold. Controls receive pseudo-event windows drawn from the empirical distribution of first-film years and are matched on baseline excess attention, baseline popularity, age, chart history, and artist-level Billboard success.

The pre-to-post increase for the selected film-linked songs exceeds that of matched controls by 5.04 excess-attention points (integrated 95\% CI [2.97, 6.87]). This interval incorporates song-level resampling, snapshot-specific decay-curve refitting, pseudo-event assignment, and rematching. All 500 pseudo-event assignment-specific contrasts are positive, with a central assignment interval of [4.24, 6.00]. This interval describes stability across assignments, rather than sampling uncertainty. The probability of exceeding the fixed pre-film tercile threshold is 39.5\% for film-linked songs and 21.5\% for controls, a difference of 18.0 percentage points (integrated 95\% CI [6.7, 27.6]). These results support an observational reactivation signature within the lower-tercile group. They do not establish that its adjusted increase is larger than that of the other terciles.

Four larger-sample definitions based on raw popularity or regression residuals yield matched differences in change ranging from 3.76 to 7.72 Spotify points, with bootstrap intervals excluding zero. Their selection criteria differ from the primary excess-attention definition. Full definition checks and integrated-inference diagnostics are reported in the SI Appendix, Pre-Film Attention and Integrated Inference. As an interpretive translation, the primary excess-attention contrast corresponds to 6.88 decay-equivalent years (integrated 95\% CI [3.89, 10.22]). This is the age difference along the fitted curve corresponding to the attention contrast, not an additional result about years of memory recovered.

\begin{figure}[!htbp]
\centering
\includegraphics[width=0.9\linewidth]{Figures/Main/figure2_expected_memory_excess_dormant.pdf}
\caption{Expected-attention baselines and temporal renewal. Panel A: snapshot-specific collective-memory decay fits (dark lines) to age means for songs with no observed film appearance (blue points). Film-linked age means (red points) are plotted for comparison. Each curve is a cross-sectional age--attention profile. Panel B: excess-attention contrasts spanning broad levels, exposure restricted to films released by each snapshot, prior-attention matching over 12,095 pairs, temporal comparisons over 93 pairs, and the primary lower-tercile difference in change. Excess attention is observed Spotify popularity minus its age/snapshot-specific fitted value. Bars are 95\% CIs, using HC1 for the broad level comparisons, a replacement-aware two-way cluster bootstrap for prior-attention matching, matched-pair standard errors for the temporal comparisons, and an integrated bootstrap for the lower-tercile comparison. Panel C: decay-equivalent translations of the same matched and temporal contrasts, shown as median values across snapshots at ages 1--58. These translations do not measure memory recovered. Panel D: distribution of 500 pseudo-event assignment-specific lower-tercile contrasts, with the integrated estimate and its 95\% CI ($n = 119$ matched treated songs per assignment). The distribution describes assignment stability, not sampling uncertainty. Fits and translations appear in the SI Appendix section \textit{Expected-Attention Baselines and MemoryDecay Diagnostics}. Selection definitions and integrated inference appear in \textit{Pre-Film Attention and Integrated Inference}.}
\label{fig:excess-memory}
\end{figure}


\begin{table}[!htbp]
\centering
\small
\caption{Excess-attention differences and temporal changes.}
\label{tab:memory-decay-key-results}
\renewcommand{\arraystretch}{1.15}
\begin{tabularx}{\linewidth}{@{}>{\raggedright\arraybackslash}p{0.175\linewidth}>{\raggedright\arraybackslash}p{0.15\linewidth}>{\raggedright\arraybackslash}p{0.155\linewidth}>{\centering\arraybackslash}p{0.145\linewidth}>{\raggedright\arraybackslash}X@{}}
\toprule
Result & Estimate & \makecell[l]{95\% CI /\\ uncertainty} & \makecell[b]{Median\\ decay-equivalent\\ years} & Interpretation \\
\midrule
Pooled film-linked excess premium & 19.12 excess-attention points & [18.86, 19.38] & -- & broad excess-attention difference \\
Film-embedded by snapshot (time-respecting) & 19.15 excess-attention points & [18.88, 19.41] & -- & time-respecting excess-attention difference \\
Prior-attention matched excess premium & 2.52 excess-attention points & [2.16, 2.87] two-way cluster bootstrap & 3.78 & matched observational comparison \\
Strict temporal event contrast & 2.45 excess-attention points & [0.62, 4.28] & 3.41 & conservative temporal check \\
Strict temporal post contrast & 2.98 excess-attention points & [1.17, 4.80] & 4.11 & conservative temporal check \\
Lower pre-film excess-attention contrast & 5.04 excess-attention points & [2.97, 6.87] integrated bootstrap & 6.88 [3.89, 10.22] & placebo-adjusted difference in change among songs in the lower pre-film excess-attention tercile \\
Exceeding the pre-film threshold & 18.0 percentage-point difference (39.5\% vs 21.5\%) & [6.7, 27.6] pp integrated bootstrap & -- & difference in above-threshold probability \\
\bottomrule
\end{tabularx}
\begin{flushleft}\footnotesize
Excess attention is observed Spotify popularity minus the fitted age/snapshot-specific expected attention from canonical collective-memory decay curves estimated on never-film age means. Intervals: HC1 for the broad and film-embedded rows; a two-way cluster bootstrap over treated and reused control song records for the matched residual, accounting for control reuse and repeated song identities; matched-pair standard errors for the strict temporal rows ($n=93$ pairs); integrated nested-bootstrap inference (song-level resampling, decay-curve refitting, pseudo-event reassignment, and rematching) for the lower-tercile excess-attention rows. Decay-equivalent years ask how many years younger a song would need to be, under the fitted expected-attention curve, to reach the shifted attention level; they are interpretive translations, not literal measures of memory recovered or randomized causal effects.
\end{flushleft}
\end{table}




\subsection*{Film appearance counts, box office, and a complementary temporal comparison}

Secondary analyses characterize how attention varies with repeated film appearances and film-level visibility. In pooled Spotify models that count only films released by each snapshot year, songs with one appearance are 9.25 popularity points above songs with no observed appearance (95\% CI [8.73, 9.76]). The difference is 16.72 points for two or three appearances (95\% CI [16.05, 17.39]) and 24.19 points for four or more (95\% CI [23.35, 25.02]). The ordering also holds within each snapshot (Figure~\ref{fig:embedding-temporal}A). Last.fm cumulative listener reach shows corresponding log-listener premiums of 1.34, 2.06, and 2.78 using lifetime movie-count categories.

For songs linked to films released by the snapshot year with recorded worldwide box office, log1p maximum worldwide box office is positively associated with Spotify popularity: 1.15 points per log unit (95\% CI [0.95, 1.35], $N=15{,}245$ song-snapshots, 4,294 normalized song-artist identities). Adjusted contrasts across terciles of the same measure show a graded pattern (Figure~\ref{fig:embedding-temporal}B). These associations are compatible with both selective reuse of more popular music and greater attention around more visible films. Counts and box office characterize reuse and film reach, without measuring individual soundtrack exposure or song prominence. The SI Appendix, Repeated Embedding and Film Visibility, reports support and selection diagnostics.

Figure~\ref{fig:embedding-temporal}C displays the broad event-time differences reported above. A complementary temporal comparison matches songs to controls with identical pre-film Spotify popularity and similar Billboard histories, retaining 93 matched pairs. It yields contrasts of 2.43 popularity points in the event year and 2.94 points in the first post period, with randomization-inference $p$-values of 0.010 and 0.0012. On the excess-attention scale, the same pairs yield 2.45 and 2.98 points (Figure~\ref{fig:embedding-temporal}D). The two scales express the same comparison. This check supports positive temporal contrasts under exact baseline-popularity matching on a different sample from the primary lower-tercile analysis.

\begin{figure}[!htbp]
\centering
\includegraphics[width=0.94\linewidth]{Figures/Main/figure3_embedding_visibility_temporal.pdf}
\caption{Film appearance counts, box office, and temporal comparisons. Panel A: snapshot-specific adjusted Spotify differences by number of films released by the snapshot year, with HC1 95\% CIs. The text reports pooled estimates with clustering by song. Panel B: adjusted contrasts across recorded worldwide box-office terciles on the continuous model's estimation sample, with the low tercile as reference and cluster-by-song 95\% CIs. The continuous estimate is 1.15 Spotify points per log unit (95\% CI [0.95, 1.35]). Box office describes film-level visibility, without measuring soundtrack prominence. Panel C: event-time contrasts on 83,785 song-snapshot observations relative to songs with no observed film appearance and observations in the immediate pre-film year. Pre denotes at least two years before the first observed film appearance. These are differences between event-time groups, not a fixed cohort's trajectory. Panel D: temporal contrasts for 93 pairs matched on identical pre-film popularity and similar chart histories. Randomization inference evaluates these observational matched comparisons. Supporting tables appear in the SI Appendix sections \textit{Repeated Embedding and Film Visibility} and \textit{Temporal Consistency Checks}.}
\label{fig:embedding-temporal}
\end{figure}



\subsection*{Other songs by the same artists show prior advantages but inconclusive changes}

We next examine other Billboard songs by the same artists that have no observed film appearance. In the pooled Spotify panel, these songs have a 2.47-point popularity advantage over songs by artists with no film-linked Billboard record, after adjusting for song age, chart history, artist catalog size, artist-level Billboard success, superstar status, and snapshot fixed effects. The difference is positive in every Spotify snapshot, and Last.fm cumulative listener reach also shows a positive association.

As with directly linked songs, an advantage is already present before placement. Other songs by artists whose Billboard music later appears in films have a 5.31-point pre-event premium in the $\leq -2$ bin. Event and early post-event bins show no abrupt increase over this level. These comparisons document prior differences between artist groups, rather than an event-induced catalog effect.

Matching on prior Spotify attention absorbs most of the catalog premium. Historical matching yields differences of 2.89, 1.38, and 1.46 points in July 2017, August 2022, and August 2025. Adding prior-attention matching with a five-point caliper reduces them to 0.46, $-0.31$, and 0.44 points, with a pooled difference of 0.31 points. Under replacement-aware cluster-by-control inference, the August 2022 contrast is not distinguishable from zero (Figure~\ref{fig:catalog}A).

Holding these matched pairs fixed, we compare prior-to-current popularity changes, with uncertainty estimated by a two-way bootstrap over treated and reused control artist identities. Differences in change are $+0.27$ points in July 2017 (95\% CI $[-0.36, 0.89]$), $-0.48$ in August 2022 ($[-1.78, 0.71]$), and $+0.40$ in August 2025 ($[-0.22, 0.95]$). The pooled descriptive difference is $+0.25$ points ($[-0.15, 0.66]$), and leave-one-cohort-out summaries remain inconclusive (Figure~\ref{fig:catalog}B). These snapshot-to-snapshot comparisons are not aligned to individual film events. They do not establish corresponding attention gains in other Billboard songs by the same artists, and their intervals do not demonstrate the absence of such gains. Full static, dynamic, matched, and change-score results appear in the SI Appendix, Artist-Catalog Attention Analyses.

\begin{figure}[!htbp]
\centering
\includegraphics[width=\linewidth]{Figures/Main/figure4_artist_catalog_regimes.pdf}
\caption{Other Billboard songs by the same artists show attenuated attention differences and inconclusive changes. The unit is the song-snapshot for songs with no observed film appearance. Catalog exposure means that the same normalized artist or band has another Billboard song linked to a film. Panel A: historical matching (open markers, paired 95\% CIs) and prior-attention matching (filled markers, replacement-aware cluster-by-control 95\% CIs) within each cohort. Connecting lines and interval-free pooled points are descriptive. The October 2016 historical estimate has no prior-attention counterpart and appears in the SI Appendix. Panel B: differences in prior-to-current attention change on the fixed matched pair sets, with 95\% CIs from a two-way bootstrap over treated and reused control artist identities. Every artist-level interval includes zero. These observational comparisons do not establish corresponding renewal in these songs or its absence. Matching, inference, and selection diagnostics appear in the SI Appendix, Artist-Catalog Attention Analyses.}
\label{fig:catalog}
\end{figure}


\section*{Discussion}

Film-linked Billboard songs have attention advantages before observed placement, yet prior visibility does not exhaust the temporal evidence. Among songs in the lower tercile of pre-film excess attention, the increase around placement exceeds that of matched controls by 5.04 excess-attention points. Positive temporal contrasts also appear in the complementary exact-pre-attention comparison. These findings support an observational account of cross-domain reactivation in which selective reuse coexists with subsequent renewal. They answer different parts of the same problem: broad comparisons reveal which songs attract more attention, and temporal comparisons ask whether attention subsequently changes beyond the selected baseline.

The expected-attention baseline connects this distinction to research on the decay and persistence of collective attention \citep{Candia2019,CandiaUzzi2021SelectiveForgetting}. It provides an explicit age- and snapshot-specific benchmark, allowing a change in attention to be evaluated relative to the attention typical of similarly aged songs. The fitted curves summarize age differences within each snapshot. They do not directly identify transitions between communicative and cultural memory, a distinction that matters when interpreting digital traces through collective-memory theory \citep{yasseri2022collective}. Nor can they establish how an individual song would have decayed without a film appearance. Matching and pseudo-event comparisons supply the additional temporal benchmark. Decay-equivalent years translate the resulting contrasts onto the same fitted curves and add an interpretive scale rather than independent evidence of renewal.

Earlier work using these soundtrack data documented broad associations between film linkage, song popularity, repeated appearances, and film commercial performance \citep{utreras2024}. The present analyses clarify the role of prior attention in those associations. Much of the broad song-level premium contracts after prior-attention matching, and positive pre-event differences make selection observable. The remaining temporal contrasts are consistent with attention renewal after accounting for measured histories. A film may introduce a song to new listeners, prompt familiar listeners to return, or coincide with wider promotion and recommendation. The data cannot separate these pathways. Cultural-memory accounts of material preserved for later reuse offer a useful interpretation \citep{Assmann1995,Assmann2008,Assmann2011}, but recognizability, emotional meaning, and latent cultural availability are not measured here.

The evidence for other music by the same artist is less conclusive. Published evidence that a new album can draw attention to earlier albums motivates the possibility of broader discovery within an artist's repertoire \citep{HendricksSorensen2009}. In our setting, other Billboard songs by artists with film-linked music already have attention advantages before placement. Prior-attention matching attenuates those differences, and matched changes are small, heterogeneous, and statistically inconclusive. These analyses therefore do not establish corresponding renewal across those songs. They neither rule out smaller increases nor test spillover to an artist's complete repertoire. Because the focal-song and catalog analyses use different samples and temporal designs, their contrast is not a formal test that attention remains confined to the soundtrack song.

The observational design leaves time-varying selection and concurrent promotion as plausible explanations for part of the temporal contrasts. Four unevenly spaced Spotify snapshots cannot locate the onset or duration of renewed attention precisely. Spotify popularity is a platform-provided 0--100 index based mainly on play volume and recency \citep{SpotifyPopularity}, rather than a stream count, and changes in the platform may affect comparisons across dates. Last.fm adds evidence about cumulative reach but cannot supply a temporal replication. Linkage uncertainty also matters: a manual sample of 200 retained song-film links confirmed 90\%, with the remaining cases ambiguous or incorrect. Confusion between performers of the same composition can misassign film identity, appearance counts, and box office. The SI documents this uncertainty without treating the sample as a recall audit. More frequent listening measurements tied to verified placements and measured audience exposure would help distinguish short-lived publicity, sustained renewal, and concurrent attention shocks.

The contribution is a joint account of prior visibility and subsequent renewal among previously successful songs reused in film. The broad associations concern thousands of Billboard records, while the primary temporal result is conditional on 119 songs with relatively lower pre-film excess attention and adequate measurement coverage. Within that scope, renewed attention exceeds the increase observed among comparable controls. The result connects selective cultural reuse to changes in collective attention, while leaving individual memory and the causal contribution of film exposure open.

\section*{Materials and Methods}

\subsection*{Data sources, construction, and linkage}

The analysis links Billboard Hot 100 song-artist records to Spotify popularity snapshots, Last.fm listener counts, IMDb song-film appearances, and available movie metadata \citep{Billboard,Spotify,Lastfm,IMDB}. The primary unit of analysis is the Billboard song-artist record: one record per distinct song-title and performer string observed on the weekly Hot 100 chart, with chart debut, total weeks on chart, and best peak position aggregated from the weekly chart entries. Billboard Hot 100 entries define a population of previously successful cultural objects. This restriction focuses the analysis on post-success persistence, decay, and reactivation rather than on initial entry into public visibility.

Spotify attention was measured using four high-coverage current-popularity snapshots: October 2016, July 2017, August 2022, and August 2025. These snapshots provide repeated song-snapshot observations of current attention on the platform. Spotify popularity is the platform-provided 0--100 compressed popularity index and is analyzed in its original scale unless otherwise stated. All Spotify analyses use these four snapshots; per-snapshot linkage coverage is documented in the SI Appendix. Last.fm attention was measured using the single July 2017 listener snapshot available in the repository: the raw \texttt{Listeners} field is treated as cumulative listener reach and analyzed as log1p cumulative listeners. No observations are imputed.

Film appearances were identified using soundtrack data collected from Soundtrack Collector and IMDb \citep{utreras2024}. During the original data collection, Soundtrack Collector film titles were standardized to lower case with special characters removed to search for the corresponding soundtrack listings in IMDb. The current analysis uses the preserved IMDb song-film records and links them to Billboard records through exact matches on separately normalized song and artist keys. A Billboard song-artist record was classified as film-linked if it matched at least one soundtrack appearance with a recorded film year. Records without an observed match were classified as having no observed film appearance in the linked data. This comparison group does not establish that a song was never used in a film, because nonmatches may reflect incomplete soundtrack coverage, undocumented uses, or limitations of exact string linkage.

Soundtrack records were linked to movie metadata through shared movie identifiers. Movie attributes were aggregated across observed film appearances and assigned to Billboard song-artist records through exact matches on alphanumeric song and artist keys. The resulting linkage provides movie metadata for 4,682 film-linked Billboard records and worldwide box-office information for 4,272. These joins use the identifiers and strings in the preserved inputs, without additional movie-title matching. Box-office coverage is reported within the primary song-film linkage; the snapshot-specific visibility analysis uses the separate alphanumeric-key linkage described in the SI Appendix.

Worldwide box-office revenue and production budget were collected from Box Office Mojo by web scraping during construction of the film dataset. Monetary amounts are used as recorded in the preserved input, without an additional monetary adjustment in the current analysis. The continuous visibility model uses log1p maximum worldwide box office. Movie genre and other available film attributes were linked through the same movie identifiers. Analyses involving box office, budget, or genre were estimated on the corresponding metadata-supported subsets. Worldwide box office is interpreted as a proxy for film-level visibility rather than as a measure of a song's prominence within a scene, soundtrack, trailer, or marketing campaign. For snapshot-specific analyses, film counts and maximum worldwide box office are restricted to films whose recorded release year is no later than the snapshot year. Box office uses the preserved cumulative totals, which do not establish revenue accrued by each snapshot date.

OpenAI Codex assisted with manuscript editing, code development and verification, and review of validation records against the supplied soundtrack evidence. Additional details on data construction, linkage, metadata coverage, and manual validation of 200 song-film links are reported in the SI Appendix, Data Construction and Platform Coverage.

\subsection*{Attention-difference models}

For snapshot $s$ and song $i$, the main attention-difference models estimate

\begin{equation}
Y_{is}=\alpha_s+\beta\,\mathrm{FilmLinked}_i+f(\mathrm{Age}_{is})+\gamma X_i+\epsilon_{is},
\end{equation}

where $f(\mathrm{Age}_{is})$ includes age and age squared, and $X_i$ includes Billboard weeks and peak chart position. Pooled Spotify models include snapshot fixed effects; snapshot-specific models omit $\alpha_s$. Spotify outcomes are platform popularity points. Last.fm uses log1p cumulative listeners because listener counts are raw count outcomes. Figure~\ref{fig:attention-matching}A additionally reports standardized coefficients for visual comparability across Spotify and Last.fm by standardizing the outcome within each platform/snapshot before estimating the same model. Snapshot-specific regressions use HC1 robust standard errors. Pooled Spotify models with repeated song-snapshot observations are interpreted with cluster-by-song inference where used as primary repeated-observation specifications; HC1 and cluster-by-song variants are reported in the SI Appendix, Inference and Robustness Checks. These are regression-adjusted observational associations, not randomized contrasts.

\subsection*{Expected-attention baseline (canonical collective-memory decay fit)}

The expected-attention baseline is estimated only on never-film Spotify observations from the four high-coverage snapshots \citep{Candia2025DataScienceCollectiveMemory}. The baseline is specific to age and snapshot: it does not include Billboard weeks, peak position, artist-level controls, or any second-stage Billboard-history correction in the expected value. The canonical collective-memory decay function from \citet{Candia2019} is

\begin{equation}
S(t)=N\left[\exp(-(p+r)t)+\frac{r\,\{\exp(-qt)-\exp(-(p+r)t)\}}{p+r-q}\right],
\end{equation}

equivalently
\begin{equation}
S(t)=\frac{N[(p-q)\exp(-(p+r)t)+r\exp(-qt)]}{p+r-q}.
\end{equation}

To estimate the expected-attention curve, we first restricted the baseline sample to never-film Spotify observations in each high-coverage snapshot. Within each snapshot, songs were grouped by integer age in years, and we computed the mean Spotify popularity for each observed age. We then fit the canonical collective-memory decay function to these never-film age means by minimizing the discrepancy between $\log(\mathrm{observed})$ and $\log S(t)$, where $S(t)$ is the fitted expected-attention curve at age $t$. The nonlinear fit uses the full MemoryDecay grid of starting values to reduce sensitivity to local optima, enforces positivity of $N$, $p$, and $r$ and $0<q<p+r$ through a log/logit reparametrization, and gives early-age observations greater weight (inverse age, capped at the 99th percentile) because the decay curve changes most rapidly near release. Spotify popularity is the platform-provided 0--100 compressed popularity index and is not pre-logged before this log-response fit. Age $t$ is measured in years; $\log(\mathrm{age})$ is used only in the separate log-normal comparison model. Exponential and log-normal comparisons and fit diagnostics are reported in the SI Appendix, Expected-Attention Baselines and MemoryDecay Diagnostics. Last.fm is not used to estimate the expected-attention curve because it contains one cumulative listener snapshot and its young-age profile is not monotonically decaying.

\subsection*{Excess attention and decay-equivalent years}

For Spotify song $i$ in snapshot $s$, excess attention is

\begin{equation}
\mathrm{ExcessAttention}_{is}=\mathrm{ObservedSpotify}_{is}-\widehat{S}_s(\mathrm{Age}_{is}),
\end{equation}

where $\widehat{S}_s$ is the snapshot-specific fitted MemoryDecay curve. Positive excess attention means the song is above the fitted age/snapshot baseline for Billboard songs not observed in films. Decay-equivalent years translate an excess-attention contrast onto the fitted memory curve. For a song of age $t$ and contrast $\Delta$, we solve for $t^\star<t$ such that $\widehat{S}_s(t^\star)=\widehat{S}_s(t)+\Delta$ and report $t-t^\star$. These are interpretive translations of observational contrasts, not literal measures of individual or collective memory.

\subsection*{Matched observational comparisons}

Song-level matched comparisons use nearest-neighbor matching with replacement. Historical matching balances chart year, song age, Billboard weeks, peak position, and artist-level Billboard success. Prior-attention matching additionally balances prior Spotify popularity using a five-point caliper when a previous high-coverage snapshot is available. Balance is assessed using standardized mean differences. Because controls can be reused, matched uncertainty is evaluated with replacement-aware bootstrap and cluster-by-control inference where reuse affects the estimand; paired intervals are retained as simpler descriptive diagnostics. For the pooled prior-attention matched-excess comparison, the primary interval is a two-way cluster bootstrap that independently resamples treated song records and reused control song records over the fixed 12,095-pair set, accounting jointly for control reuse under matching with replacement and for repeated song identities across snapshots; naive paired and one-way cluster alternatives are reported in the SI Appendix, Matched Song-Level Comparisons. These are observational comparisons conditional on matched covariates. They do not make film placement random or exogenous.

\subsection*{Pre-film attention and placebo comparisons}

The primary pre-film attention comparison selects film-linked songs in the bottom tercile of pre-film excess attention under the MemoryDecay expected-attention baseline, among film-linked songs with an observed Spotify snapshot strictly before and at or after the first film year. Baseline attention is measured at the latest observed snapshot strictly before the first film year. This criterion identifies relatively low pre-film excess attention within the eligible group, without requiring attention below the expected curve, an observed prior decline, or a minimum song age. Billboard inclusion establishes historical chart presence, without implying a particular level of attention before film placement. The above-threshold outcome indicates whether a song exceeds the fixed pre-film tercile threshold at the first post-film snapshot. Songs are identified by their exact Billboard record identity so that no distinct records are collapsed. Reduced-form raw definitions--the bottom tercile of residuals from a regression of baseline popularity on age, age squared, Billboard weeks, and peak position; low baseline attention; historically successful low attention; and older low attention--are retained as robustness checks. Never-film controls receive pseudo-event windows drawn from the empirical first-film-year distribution and are matched to the selected film-linked songs on baseline excess attention, baseline popularity, age, chart history, and artist-level Billboard success. The estimand is treated change minus matched control change. Final inference for this comparison integrates over 500 valid pseudo-event assignments and uses a nested bootstrap (250 replications) that resamples never-film song trajectories to refit the snapshot-specific MemoryDecay curves with the same floor-age aggregation and full multi-start estimator as the point-estimate analysis, recomputes expected and excess attention, reconstructs the lower tercile within the eligible film-linked sample, resamples the selected treated records, reassigns pseudo-events to controls, and rematches controls inside each replication. Permutation tests and single-assignment robustness checks are reported in the SI Appendix, Pre-Film Attention and Integrated Inference.

\subsection*{Repeated embedding and film visibility}

Repeated film embedding is measured using movie-count categories: no observed film appearance, one film appearance, two or three film appearances, and four or more film appearances. The main Spotify repeated-embedding models use time-respecting movie-count categories that include only films released by each Spotify snapshot year and estimate cluster-by-song intervals across repeated song-snapshot observations after adjusting for age, age squared, Billboard weeks, peak chart position, and snapshot fixed effects. Last.fm repeated-embedding models use lifetime movie-count categories with log1p cumulative listeners as the outcome because only one Last.fm listener snapshot is available. The lifetime-count comparison for Spotify is reported in the SI Appendix as a robustness check on exposure timing.

Film visibility is measured using available worldwide box office after exact internal-ID linkage. The main Spotify visibility model uses log1p maximum recorded worldwide box office among films released by the snapshot year. It is estimated on song-snapshot observations with this measure available, with cluster-by-song intervals and adjustment for the film count through that year. These analyses are observational heterogeneity checks. Box office is interpreted as film-level visibility, not as a measure of scene-level soundtrack prominence. Repeated-embedding and visibility specifications are reported in the SI Appendix, Repeated Embedding and Film Visibility.

\subsection*{Temporal checks}

Broad event-time models define event bins relative to the first observed film appearance of a Billboard song-artist record and explicitly report pre-event differences. These models adjust for age, age squared, Billboard weeks, peak position, and snapshot fixed effects, and are interpreted as temporal consistency checks that also reveal selection into film placement. The strict exact-pre-attention matched comparison is retained as a conservative temporal check. It matches treated songs to never-film controls with identical baseline Spotify attention and similar Billboard histories, using October 2016 as the baseline for the 2017 cohort and July 2017 as the baseline for the 2022 cohort. It is interpreted as temporal consistency with reactivation under stronger comparability, not as causal identification. Full support, balance diagnostics, and randomization-inference checks are reported in the SI Appendix, Temporal Consistency Checks.

\subsection*{Artist-catalog models}

Artist catalogs are defined by normalized artist/band strings. The static catalog analysis is restricted to songs never linked to films. Catalog exposure equals one when the same artist has another Billboard song linked to a film; controls are songs by artists with no film-linked Billboard song. Regression models include song-level chart history and artist-level historical controls, including artist catalog size, artist-level Billboard success, and superstar status where specified. Dynamic catalog analyses define event time as the first film appearance of any other Billboard song by the same artist. Matched catalog comparisons use historical and prior-attention matching designs analogous to the song-level matched comparisons and are interpreted with replacement-aware uncertainty checks. A matched change-score comparison re-expresses the fixed prior-attention matched pairs as prior-to-current attention changes without rematching. Because catalog exposure is defined at the artist level, its primary uncertainty is a two-way bootstrap that independently resamples treated artist identities and reused control artist identities and reweights pairs by the product of the artist-cluster draw counts (2,000 deterministic replications), with naive paired and one-way artist-level alternatives reported in the SI Appendix. These comparisons describe attention levels and changes among other Billboard songs by the same artists. The matched snapshot-to-snapshot comparisons are not aligned to individual film events and do not identify causal spillover. Full static, dynamic, matched, and change-score catalog outputs are reported in the SI Appendix, Artist-Catalog Attention Analyses.


\section*{Acknowledgments}

C.C. acknowledges financial support from ANID FONDECYT/REGULAR 1262601. The authors also acknowledge the support of the Instituto de Data Science at Universidad del Desarrollo.


\section*{Author Contributions}

C.C. conceived the study and designed the research. C.C. developed the analytical strategy, performed the analyses, interpreted the results, prepared the figures and tables, and wrote and revised the manuscript. C.U. collected the original soundtrack and film data, documented data construction and linkage methods, and manually validated song-film links. F.D. contributed to the literature review and manuscript drafting and revision.


\section*{Competing Interests}

The authors declare no competing interests.

\section*{Data, Materials, and Software Availability}
Analysis code, data dictionaries, and processed data that can be redistributed will be deposited in the Open Science Framework upon publication. Third-party inputs will be shared where redistribution is permitted; source descriptions and reconstruction instructions will accompany inputs that cannot be redistributed.


\bibliographystyle{unsrtnat}
\bibliography{bibliografia}

\end{document}
